\chapter{Purity, Verdier duality, and Poincare duality}
\label{ch:duality}

Duality supplies the trace map, reciprocal Frobenius eigenvalues, and the
self-dual parabolic cohomology used in the curve estimate.  We prove the
purity statement needed in this book for regular schemes of finite type over
a perfect field.  Gabber's more general absolute-purity theorem is recorded
separately and is not claimed as a consequence of the proof capsule below.

\section{Cohomology with supports and purity}

Let \(i:Z\hookrightarrow X\) be a closed immersion, with open complement
\(j:U\hookrightarrow X\).  The exceptional inverse image is characterized by
\[
 R\Gamma_Z(X,K)=R\Gamma(Z,Ri^!K),
\]
and localization gives
\[
 i_*Ri^!K\longrightarrow K\longrightarrow Rj_*j^*K\longrightarrow.
\tag{7.1}\label{eq:costalk-localization}
\]

\begin{theorem}[Purity in the geometric case]
\label{thm:absolute-purity}
Let \(k\) be a perfect field of characteristic different from \(\ell\), and
let \(i:Z\hookrightarrow X\) be a regular closed immersion of pure
codimension \(c\), where \(Z\) and \(X\) are regular schemes of finite type
over \(k\).  Equivalently, \(Z\) and \(X\) are smooth over \(k\).  Then the
fundamental class induces a canonical isomorphism
\[
 Ri^!\Lambda_n\xrightarrow{\sim}\Lambda_n(-c)[-2c].
\tag{7.2}\label{eq:absolute-purity}
\]
It is compatible with etale base change, transverse composition of regular
immersions, cup products, and reduction from \(\Lambda_n\) to
\(\Lambda_{n-1}\).

If \(f:X\to Y\) is smooth of pure relative dimension \(d\), then
\[
 Rf^!L\xrightarrow{\sim}f^*L(d)[2d]
\tag{7.3}\label{eq:smooth-purity}
\]
for \(L\in D^b_c(Y,\Lambda_n)\).  The same formulas hold for
\(\mathbf Z_\ell\)- and \(\mathbf Q_\ell\)-coefficients.
\end{theorem}

\begin{proof}[Proof capsule]
For the zero section \(i:\{0\}\hookrightarrow\mathbf A^1_k\), the localization
sequence and Kummer sequence calculate
\[
 H^r_{\{0\}}(\mathbf A^1_k,\Lambda_n)=
 \begin{cases}
 \Lambda_n(-1),&r=2,\\
 0,&r\ne2.
 \end{cases}
\tag{7.4}\label{eq:purity-line}
\]
The boundary of the Kummer class of the coordinate is the generator.
Multiplying the coordinate by a unit does not change its supported boundary.

Kunneth gives the corresponding statement for the zero section of
\(\mathbf A^c\).  Etale descent globalizes the product of the \(c\) Kummer
classes to the Thom class of a rank-\(c\) vector bundle.  For a general
regular immersion form the deformation to the normal cone
\[
 D_ZX=
 \operatorname{Bl}_{Z\times\{0\}}(X\times\mathbf A^1)
 \setminus \operatorname{Bl}_{Z}(X).
\]
Its fiber over \(1\) is \(X\), and its fiber over \(0\) is the normal bundle
\(N_{Z/X}\).  Smoothness of the pair over the perfect field gives the
specialization isomorphism for cohomology with supports.  Under this
isomorphism, the fundamental class of \(Z\subset X\) becomes the Thom class
of the zero section \(Z\subset N_{Z/X}\).  The vector-bundle calculation
therefore proves \eqref{eq:absolute-purity}.  Products of Kummer classes prove
composition compatibility, and base change of the deformation space proves
base-change compatibility.

For a smooth \(f\), the diagonal
\(\Delta_f:X\hookrightarrow X\times_YX\) is a regular immersion of
codimension \(d\).  Apply purity to the diagonal and the two projection
adjunctions.  This identifies \(Rf^!\Lambda_n\) with
\(\Lambda_n(d)[2d]\); the projection formula gives
\eqref{eq:smooth-purity}.  Compatibility in \(n\), normalization, and
inversion of \(\ell\) give the adic assertions.
\end{proof}

\begin{remark}[Gabber's absolute purity]
For a regular closed immersion between excellent regular schemes on which
\(\ell\) is invertible, the same isomorphism
\(Ri^!\Lambda_n\simeq\Lambda_n(-c)[-2c]\) is Gabber's absolute-purity theorem
\cite{SGA4half}.  Its mixed-characteristic proof is genuinely
deeper than deformation to the normal cone and is not supplied here.  The
geometric case above is sufficient for all varieties, smooth morphisms, and
Lefschetz pencils used in this book.
\end{remark}

\begin{corollary}[Gysin morphism]
\label{cor:gysin}
Under the hypotheses of Theorem~\ref{thm:absolute-purity}, adjunction gives
\[
 i_*:H^r(Z,\Lambda_n(s))
 \longrightarrow H^{r+2c}_Z(X,\Lambda_n(s+c)).
\]
Gysin morphisms compose and commute with transverse pullback.
\end{corollary}

\section{Exceptional inverse image and Verdier duality}

\begin{theorem}[Existence of \(Rf^!\)]
\label{thm:exceptional-pullback}
For a separated finite-type morphism \(f:X\to Y\), the functor
\[
 Rf_!:D^b_c(X,\Lambda_n)\longrightarrow D^b_c(Y,\Lambda_n)
\]
has a right adjoint \(Rf^!\), characterized by
\[
 R\operatorname{Hom}_Y(Rf_!K,L)
 \simeq R\operatorname{Hom}_X(K,Rf^!L).
\tag{7.5}\label{eq:fshriek-adjunction}
\]
It is compatible with composition, equals \(f^*\) for etale \(f\), satisfies
\eqref{eq:smooth-purity} for smooth \(f\), and agrees with supported costalks
for a closed immersion.
\end{theorem}

\begin{proof}[Proof capsule]
For a compactification \(f=\bar f j\), set
\(Rf^!=j^*R\bar f^!\).  The right adjoint for a proper map is constructed
first for finite maps by sheaf Hom, for projective space by the trace class,
for projective maps after a closed embedding, and for proper maps by
cohomological descent from projective envelopes.  The adjunction is checked
at each stage, and the uniform cohomological bound preserves the bounded
constructible category.  A common refinement of compactifications and
uniqueness of right adjoints identify the results.  The same uniqueness gives
composition and all coherence isomorphisms; see
\cite[Expos\'e XVIII]{SGA4}.
\end{proof}

Let \(a_X:X\to\operatorname{Spec}k\), with \(k\) separably closed, and define
\[
 \omega_X:=Ra_X^!\mathbf Q_\ell,\qquad
 D_XK:=R\mathcal Hom(K,\omega_X).
\]
If \(X\) is smooth of pure dimension \(d\), then
\[
 \omega_X\simeq\mathbf Q_\ell(d)[2d].
\tag{7.6}\label{eq:smooth-dualizing}
\]

\begin{theorem}[Verdier duality]
\label{thm:verdier-duality}
For \(K\in D^b_c(X,\mathbf Q_\ell)\):
\begin{enumerate}
\item \(K\to D_XD_XK\) is an isomorphism;
\item for separated finite-type \(f:X\to Y\),
\[
 D_YRf_!K\xrightarrow{\sim}Rf_*D_XK;
\tag{7.7}\label{eq:verdier-push}
\]
\item for \(L\in D^b_c(Y,\mathbf Q_\ell)\),
\[
 D_Xf^*L\xrightarrow{\sim}Rf^!D_YL;
\tag{7.8}\label{eq:verdier-pull}
\]
\item these maps commute with composition, base change, and Weil structures.
\end{enumerate}
\end{theorem}

\begin{proof}
Biduality is etale local and additive in localization triangles.  On a smooth
stratum where \(K\) is a bounded complex of finite free lisse sheaves, it is
ordinary biduality for finite-dimensional vector spaces, shifted and twisted
by \eqref{eq:smooth-dualizing}.  Noetherian induction over a finite
stratification proves the global assertion.

To prove \eqref{eq:verdier-push}, test both sides against an arbitrary complex
and use, successively, tensor--Hom adjunction, the projection formula, and
\eqref{eq:fshriek-adjunction}; Yoneda gives the isomorphism.
Equation \eqref{eq:verdier-pull} is its adjoint form.  The maps are uniquely
specified by units, counits, and evaluation, so their coherence and Weil
equivariance follow from the corresponding properties of the six operations.
\end{proof}

\section{Trace and Poincare duality}

For a smooth separated \(d\)-fold \(X/k\), the counit
\(Ra_{X!}Ra_X^!\mathbf Q_\ell\to\mathbf Q_\ell\) is the trace map
\[
 \operatorname{tr}_X:
 H_c^{2d}(X,\mathbf Q_\ell(d))\longrightarrow\mathbf Q_\ell.
\tag{7.9}\label{eq:trace-map}
\]
It is normalized so that the product of the Kummer fundamental classes on
\(\mathbf A^d\) maps to \(1\).

\begin{theorem}[Poincare duality]
\label{thm:poincare-duality}
Let \(X\) be smooth, separated, and pure of dimension \(d\) over a separably
closed field, and let \(L\) be a lisse \(\mathbf Q_\ell\)-sheaf.  Cup product,
evaluation, and trace give a perfect pairing
\[
 H_c^r(X,L)\times H^{2d-r}(X,L^\vee(d))
 \longrightarrow\mathbf Q_\ell.
\tag{7.10}\label{eq:poincare-pairing}
\]
Equivalently,
\[
 R\Gamma_c(X,L)^\vee
 \simeq R\Gamma(X,L^\vee(d))[2d].
\tag{7.11}\label{eq:poincare-derived}
\]
The pairing is functorial and Weil equivariant.
\end{theorem}

\begin{proof}
Apply \eqref{eq:verdier-push} to \(a_X\):
\[
 R\operatorname{Hom}(R\Gamma_c(X,L),\mathbf Q_\ell)
 \simeq R\Gamma(X,D_XL).
\]
Smooth purity identifies \(D_XL\) with \(L^\vee(d)[2d]\).
Finite-dimensionality makes the resulting cohomological pairing perfect.
The construction uses only Weil-equivariant evaluation and adjunction maps.
\end{proof}

\begin{proposition}[Frobenius reciprocity]
\label{prop:frob-duality}
Let \(X_0/\mathbf F_q\) be smooth proper of dimension \(d\).  The untwisted
pairing
\[
 H^r(X,L)\times H^{2d-r}(X,L^\vee)
 \longrightarrow\mathbf Q_\ell(-d)
\tag{7.12}\label{eq:untwisted-pairing}
\]
satisfies
\[
 \langle Fv,Fw\rangle=q^d\langle v,w\rangle.
\tag{7.13}\label{eq:frob-similitude}
\]
Thus an eigenvalue \(\alpha\) on \(H^r(X,L)\) corresponds, with the same
algebraic multiplicity, to \(q^d/\alpha\) on
\(H^{2d-r}(X,L^\vee)\).
\end{proposition}

\begin{proof}
The target is \(\mathbf Q_\ell(-d)\), on which geometric Frobenius acts by
\(q^d\).  Perfectness identifies the operator on the second space with
\(q^d(F^{-1})^{\mathsf t}\), proving both assertions.
\end{proof}

\begin{corollary}[Relative Poincare duality]
\label{cor:relative-poincare}
If \(f:X_0\to Y_0\) is smooth proper of relative dimension \(d\), then
\[
 (R^rf_*L)^\vee\simeq R^{2d-r}f_*L^\vee(d)
\]
as Weil sheaves.
\end{corollary}

\begin{proof}
Apply Poincare duality to every geometric fiber.  Proper and smooth base
change identify the terms with stalks of the displayed sheaves; conservativity
of geometric stalks glues the isomorphism.
\end{proof}

\section{Local duality and middle extension on a curve}

Let \(C\) be a smooth proper geometrically connected curve, let
\(j:U\hookrightarrow C\) be nonempty open, and let \(L\) be lisse on \(U\).
At \(s\in C\setminus U\), write \(I_s\) for inertia and
\(V=L_{\bar\eta}\).  Then
\[
 (j_*L)_{\bar s}=V^{I_s},\qquad
 (R^1j_*L)_{\bar s}=H^1(I_s,V),\qquad
 R^rj_*L=0\quad(r>1).
\tag{7.14}\label{eq:curve-local-stalks}
\]
Indeed, wild inertia is pro-\(p\), while the maximal pro-\(\ell\) quotient of
tame inertia has \(\ell\)-cohomological dimension one.

The intermediate extension in the perverse heart is
\[
 j_{!*}L[1]
 :=\operatorname{im}({}^pj_!L[1]\to{}^pj_*L[1]).
\]
On a smooth curve the stalk calculation above gives
\[
 j_{!*}L[1]\simeq(j_*L)[1].
\tag{7.15}\label{eq:curve-middle-extension}
\]

\begin{theorem}[\(j_*\)-local duality on curves]
\label{thm:curve-middle-duality}
There is a canonical Weil-equivariant isomorphism
\[
 D_C((j_*L)[1])
 \xrightarrow{\sim}(j_*L^\vee(1))[1].
\tag{7.16}\label{eq:jstar-local-duality}
\]
Consequently,
\[
 H^1(C,j_*L)\times H^1(C,j_*L^\vee(1))
 \longrightarrow\mathbf Q_\ell
\tag{7.17}\label{eq:curve-middle-pairing}
\]
is perfect, and
\[
 H^1(C,j_*L)
 =
 \operatorname{im}\bigl(H_c^1(U,L)\to H^1(U,L)\bigr).
\tag{7.18}\label{eq:parabolic-cohomology}
\]
\end{theorem}

\begin{proof}
Verdier duality exchanges \(Rj_!\) and \(Rj_*\) and is \(t\)-exact for the
middle-perverse \(t\)-structure.  Since
\(D_U(L[1])=L^\vee(1)[1]\), it takes the image defining \(j_{!*}\) to the
opposite image and gives
\[
 D_C(j_{!*}L[1])\simeq j_{!*}L^\vee(1)[1].
\]
Use \eqref{eq:curve-middle-extension} on both sides.  The local content is
Tate duality for inertia: evaluation and the residue map pair the invariant
boundary term on one side with the coinvariant boundary term on the other.

Global Verdier duality now gives \eqref{eq:curve-middle-pairing}.  Finally,
the localization long exact sequences show that both
\(H^1(C,j_*L)\) and the image of compactly supported cohomology are the kernel
of
\[
 H^1(U,L)\longrightarrow
 \bigoplus_{s\in C\setminus U}H^1(I_s,V).
\]
This proves \eqref{eq:parabolic-cohomology}.
\end{proof}

\section{Formalization checklist}

\begin{enumerate}
\item State geometric purity with its actual scope: regular finite-type pairs
over a perfect field.  Record Gabber's excellent-regular theorem as a separate
external result, not as a consequence of Popescu approximation.
\item Construct fundamental and Thom classes, deformation to the normal cone,
and their base-change and composition laws.
\item Export \(Rf^!\) with its adjunction, dualizing complexes, biduality, and
the two Verdier exchange isomorphisms.
\item Normalize the trace on affine space and test every shift and twist:
\[
 D_X(\mathbf Q_\ell)=\mathbf Q_\ell(d)[2d],\quad
 H^{2d}(X,\mathbf Q_\ell)=\mathbf Q_\ell(-d).
\]
\item Store both the twisted invariant Poincare pairing and the untwisted
\(q^d\)-similitude pairing; derive
\(\alpha\leftrightarrow q^d/\alpha\).
\item For curves, separate inertia stalk calculations, perverse intermediate
extension, local Tate duality, and the parabolic image theorem.
\end{enumerate}

